% =============================================================
%  第八单元  乘除法的应用（一）
% =============================================================
\Chapter{第八关\quad 算账的本事}
\noindent{\small\kaishu\color{tipc}（课本第八单元：乘除法的应用（一），P87--P93）}

\begin{blunder}
超市结账，收银员说“一共 78”。爸爸扫了码就走，回家对单子才发现：
多扫了一件 9 块的东西，是爸爸自己放进袋子的。
没人提醒，是爸爸自己对着小票一项一项点出来的。
——从那以后，爸爸买东西必对小票。
\end{blunder}

\begin{gamble}{课前协商：二十块预算}
\begin{itemize}
  \item 给小豆角 50 块，让小豆角在超市自己买三样东西，\textbf{自己算够不够}。
  \item 超支了怎么办、省下来的归谁，现在就谈定。
  \item 第二次玩的时候加一条：\textbf{必须买一样“给别人的”}（给妈妈、给爷爷）。
        小豆角会开始算“留给别人的那部分”，这比任何应用题都难。
\end{itemize}
\negolines
\end{gamble}

\begin{bonus}
\begin{itemize}
  \item \textbf{当一天家庭采购员}：给小豆角 50 块，全权负责买今天的菜或水果。
        \textbf{爸爸全程不插嘴}，只在最后问一句“还剩几块？为什么剩这么多？”
  \item \textbf{猜总价}：结账前两人各估一个总价，谁接近谁赢，赢的人记一颗星。
  \item \textbf{36 元挑战}：给小豆角 36 块，让小豆角用“同一种商品”正好花完
        （可以买单价 3、4、6、9、12 的东西）。这题能玩很久。
\end{itemize}
\end{bonus}

\begin{unitinfo}
\textbf{课本 P87--P93}\quad \textbf{4 次课，每次 20 分钟} \\[4pt]
\textbf{这一关到底在干嘛}：前面学的乘和除，这一章全拿来解决真问题。\\
\textbf{逻辑是这么走的}：\ 先判断“知道什么、求什么” $\to$ 判断用乘还是除
$\to$ 想口诀 $\to$ 写算式加单位加答语 $\to$ 用乘法验一遍。\\
\textbf{一个新增难点}：“括号里最大能填几”（$4\times(\ )<18$）。
这个不用新知识，就是\textbf{试乘}——试到刚好不超过为止。
\end{unitinfo}

\subsection{开局先聊个事：这些就是生活里的算账}

长颈鹿造房子，一间用 7 块木板，63 块能造几间？
小鸟住房，6 只住一间，42 只要几间？
——这就是\textbf{知道总数，求能分几份}，用除法。

柿子 54 个，每盒装 6 个，装几盒？
——还是一样。

反过来，5 盒能装多少个？——这是\textbf{知道一份和份数，求总数}，用乘法。

\textbf{生活里算账，来回就这两个方向。}判断对了，剩下就是口诀的事。

\subsection{任务 1：先判断，再动笔}

\begin{lesson}{1}{20}
\textbf{今天要拿下的}：一看到题，先能判断用乘还是用除，说得出理由。

\begin{materials}
纸、笔、小棒（不确定时可以摆）。
\end{materials}

\begin{enumerate}
  \item \textbf{两道题一起出}（课本 P87 长颈鹿与小鸟）：
        \begin{itemize}
          \item \textbf{A}：一间房要 7 块木板，一共 63 块，能造几间？
                $\to 63\div7=9$（间），口诀“七九六十三”
          \item \textbf{B}：42 只小鸟，每 6 只住一间，要几间？
                $\to 42\div6=7$（间），口诀“六七四十二”
        \end{itemize}
        \textbf{两题都是除法}：知道总数，求“能分成几份”。
  \item \textbf{追问（课本 P88）}：造了 9 间，住了 7 间，还有几间空房？
        $9-7=2$（间）。还能住几只？$2\times6=12$（只）。
        \textbf{一道题里乘、除、减都用上了}——这就是“综合”。
  \item \textbf{每题旁边写口诀}。\textbf{写不出口诀，说明是猜的。}
  \item \textbf{只判断，不算（30 秒六道）}：
        \begin{enumerate}[label=（\arabic*）,itemsep=2pt]
          \item 每盒 8 个，5 盒一共多少个？（乘）
          \item 40 个苹果，每盒 8 个，装几盒？（除）
          \item 40 个苹果，平均装 5 盒，每盒几个？（除）
          \item 一班 35 人，7 人一组，分几组？（除）
          \item 每排 7 人，5 排一共多少人？（乘）
          \item 一支笔 6 元，5 支多少元？（乘）
        \end{enumerate}
  \item \textbf{四步法再走一遍}：圈数字 $\to$ 说关系 $\to$ 想口诀 $\to$
        写算式＋单位＋答语。
\end{enumerate}

\begin{saybox}
“这题知道总数吗？知道每份几个吗？那求什么？”\\
“小豆角用的是哪句口诀？”\\
“算出来是‘间’还是‘只’？念一遍问题再写单位。”
\end{saybox}

\begin{passbox}
\begin{itemize}
  \item 六道判断题全对。
  \item 两道应用题的算式、口诀、单位、答语都齐。
\end{itemize}
\end{passbox}
\end{lesson}

\subsection{任务 2：括号里最大能填几}

\begin{lesson}{2}{20}
\textbf{今天要拿下的}：会用“试乘”解决 $4\times(\ )<18$，并且知道填完要验。

\begin{materials}
小棒 18 根、纸、笔。
\end{materials}

\begin{enumerate}
  \item \textbf{给情境}（课本 P88）：18 只小松鼠分组，每组只数一样。
        每组 4 只，\textbf{最多}分几组？
  \item \textbf{先摆}：18 根小棒，4 根一堆。堆 4 堆用掉 16 根，还多 2 根——
        \textbf{不够再堆一堆}。所以最多 4 组。
  \item \textbf{写成算式}：$4\times(\ )<18$，括号里最大填几？
        \begin{itemize}
          \item 试乘：$4\times4=16$（小于 18，行）；$4\times5=20$（超过 18，不行）
          \item 所以最大填 \textbf{4}。
        \end{itemize}
  \item \textbf{第二种办法（用除法）}：$18\div4$——想“四五二十”太大、
        “四四十六”合适，商 4，剩 2。所以填 4。
        \textbf{两种办法都教，小豆角挑顺手的用。}
  \item \textbf{连做八道}（课本 P89）：
        \begin{center}
        \small
        \begin{tabular}{@{}llll@{}}
        $8\times(\ )<70$ & $6\times(\ )<29$ & $9\times(\ )<50$ & $(\ )\times9<55$ \\
        $(\ )\times7<65$ & $(\ )\times5<48$ & $5\times(\ )<36$ & $8\times(\ )<39$ \\
        \end{tabular}
        \end{center}
        答案：8、4、5、6、9、9、7、4。
  \item \textbf{易错对比}：$8\times(\ )<70$ 为什么不是 9？
        （$8\times9=72$，比 70 大，不行）——\textbf{填完必须乘一遍验}。
  \item \textbf{小兔安家}（课本 P89）：24 只小兔。
        \begin{center}
        \small
        \begin{tabular}{@{}lcccc@{}}
        \toprule
        每间住几只 & 4 & 6 & 3 & 8 \\
        \midrule
        需要几间 & \blank[1cm] & \blank[1cm] & \blank[1cm] & \blank[1cm] \\
        \bottomrule
        \end{tabular}
        \end{center}
        答案：6、4、8、3。\textbf{顺便发现}：每间住得越多，需要的间数越少。
\end{enumerate}

\begin{saybox}
“$4\times5=20$，20 比 18 大还是小？能填 5 吗？”\\
“除了一个个试，还能怎么想？”（用除法：$18\div4$）\\
“填完小豆角怎么确认它对？”（乘一遍，看有没有超过右边）
\end{saybox}

\begin{pitbox}
\textbf{最常见的错}：填了一个让乘积\textbf{超过}右边数的值
（比如 $8\times(\ )<70$ 填 9）。\\
\textbf{解法}：\textbf{这类题填完必须乘一遍}。
养成“填完就验”，比多做十道题有用。
\end{pitbox}

\begin{passbox}
\begin{itemize}
  \item 八道全对，\textbf{每道都有验算痕迹}。
  \item 能用两种办法（试乘 / 用除法）解同一道。
\end{itemize}
\end{passbox}
\end{lesson}

\subsection{任务 3：一页图里全是应用}

\begin{lesson}{3}{20}
\textbf{今天要拿下的}：同一幅图里把乘、除、倍、加减都用一遍。

\begin{materials}
课本 P90 农家小院、笔。
\end{materials}

\begin{enumerate}
  \item \textbf{先读图说话}：让小豆角把图里能看到的数学信息\textbf{全说出来}：
        柿子 54 个、每盒装 6 个、墙上的玉米、兔子、豇豆 42 厘米、扁豆 7 厘米……
        \textbf{这一步叫“读图”。考试里八成的错，是因为信息没读全就开始写。}
  \item \textbf{逐题解决}（每题都走四步法）：
        \begin{enumerate}[label=（\arabic*）,itemsep=3pt]
          \item 54 个柿子，每盒 6 个，装几盒？$54\div6=9$（盒）
          \item 豇豆 42 厘米，扁豆 7 厘米，豇豆是扁豆的几倍？$42\div7=6$（倍）
          \item 平均每个笼子住几只兔子？（让小豆角自己数出两个数再列式）
          \item 可以装几筐？（看每筐装几个）
        \end{enumerate}
  \item \textbf{接着追问}（课本 P91）：
        \begin{itemize}
          \item 多少个柿子刚好装满 5 盒？$6\times5=30$（个）——\textbf{这是乘法}
          \item 已经摘了 30 个，树上还剩多少？$54-30=24$（个）
          \item 剩下 24 个还能装几盒？$24\div6=4$（盒）
        \end{itemize}
        \textbf{一道题，乘、减、除连着用。}这就是这一章的核心。
  \item \textbf{自己提问题}：让小豆角看图提两个用乘法的、两个用除法的，写下来并解答。
  \item \textbf{36 元能买什么}（课本 P91）：给 3、4、5、6、7、9 元几种商品，
        “李老师买了同一种商品正好花 36 元，可能买了什么？买了多少？”
        \textbf{答案不止一个}：9 元的买 4 个、6 元的买 6 个、4 元的买 9 个、
        3 元的买 12 个。\textbf{让小豆角尽量多找几种}，这是很好的发散练习。
\end{enumerate}

\begin{saybox}
“图里一共有哪些数字？都找全了吗？”\\
“这道为什么用乘法不用除法？”\\
“36 元还能买什么？再想一种。”\\
“单位写了吗？答语写了吗？”
\end{saybox}

\begin{passbox}
\begin{itemize}
  \item 能说全图里的信息（至少六条）。
  \item 四道题算式对、单位对、答语完整。
  \item “36 元”那题能找出三种以上买法。
\end{itemize}
\end{passbox}
\end{lesson}

\subsection{任务 4：自己编题 + 综合过关}

\begin{lesson}{4}{20}
\textbf{今天要拿下的}：给算式能编出生活问题；一道综合题从头做到尾。

\begin{materials}
纸、笔、错题本。
\end{materials}

\begin{enumerate}
  \item \textbf{给算式编事}（每次 1 分钟，编四个）：
        $7\times8=56$、\ $48\div6=8$、\ $36\div9=4$、\ $9\times7=63$。
        \textbf{要求人物、东西、单位都对得上。}
        （“48 颗糖，每 6 颗装一袋，能装 8 袋。”）
  \item \textbf{反过来}：爸爸说一个生活情境，小豆角写算式（五个）。
  \item \textbf{综合过关题}：
        \begin{tipbox}
        学校要给 6 个班每班发 7 本图书。\\
        （1）一共需要多少本？\\
        （2）图书角现在有 40 本，够吗？还差几本？\\
        （3）学校又买来 8 本，现在一共 48 本，平均分给 8 个小组，每组几本？
        \end{tipbox}
        答案：（1）$6\times7=42$（本）；
        （2）$42-40=2$（本），\textbf{不够，还差 2 本}；
        （3）$48\div8=6$（本）。
  \item \textbf{整理错题}：把本单元和前面各单元快检的错题看一遍，
        挑出\textbf{同一个错因出现三次以上}的，那就是期末重点补的地方。
\end{enumerate}

\begin{passbox}
\begin{itemize}
  \item 四个算式都能编出合理的故事（人、物、单位都对）。
  \item 综合题三步全对。
\end{itemize}
\end{passbox}
\end{lesson}

\subsection{翻车现场（对着看）}

\begin{center}
\small
\begin{tabularx}{\linewidth}{@{}p{3.8cm}p{4.0cm}X@{}}
\toprule
\textbf{小豆角写出来的样子} & \textbf{其实是想错了什么} & \textbf{爸爸回一句什么} \\
\midrule
$8\times(\ )<70$ 填 9 & 没验算 & 填完必须乘一遍。 \\
\addlinespace[2pt]
不知道用乘还是除 & 没先判断“知道什么、求什么” & 走四步法：圈 → 说 → 想 → 写。 \\
\addlinespace[2pt]
单位写“个”其实该写“盒” & 没回头读问题 & 写完念一遍问的是什么。 \\
\addlinespace[2pt]
答语只写个数字 & 书写习惯 & “答：……。”一整句写出来。 \\
\addlinespace[2pt]
除法不验算 & 没检查习惯 & 商 $\times$ 除数 $=$ 被除数。 \\
\addlinespace[2pt]
图里的信息没读全 & 看图太快 & 先“说出来”，再动笔。 \\
\bottomrule
\end{tabularx}
\end{center}

\subsection{抽查：随机开枪（10 题）}

\begin{enumerate}[label=\arabic*.,itemsep=4pt]
  \item $63\div7=$ \blank[1.4cm] \qquad 2. $42\div6=$ \blank[1.4cm] \qquad
  3. $54\div6=$ \blank[1.4cm]
  \item $8\times(\ )<70$，最大填 \blank[1.4cm]。
  \item $9\times(\ )<50$，最大填 \blank[1.4cm]。
  \item 每间房住 6 只小鸟，42 只要 \blank[1.4cm] 间。
  \item 63 块木板，每间用 7 块，能造 \blank[1.4cm] 间。
  \item 豇豆 42 厘米，扁豆 7 厘米，豇豆是扁豆的 \blank[1.4cm] 倍。
  \item 每盒 6 个，9 盒一共 \blank[1.4cm] 个。
  \item 54 个柿子，每盒 6 个，装 \blank[1.4cm] 盒。
\end{enumerate}

\begin{ansbox}
1. 9\quad 2. 7\quad 3. 9\quad 4. 8\quad 5. 5\quad 6. 7\quad
7. 9\quad 8. 6\quad 9. 54\quad 10. 9 \\
\textbf{全对} $\to$ 可以进期末复习。\textbf{错两题以上} $\to$ 回对应的课补 20 分钟。
\end{ansbox}
